% GENERATED FILE. Edit canonical lessons, then run npm run generate:books.
% canonical-source-hash: fnv1a64:e3b68577f24fbfcf
% canonical-lessons: KA-C380-manavolisu, KA-C380-prerepisu, KA-C380-manadattu-madu, KA-C380-amisa, KA-C380-akarsisu

\chapter{Winning People Over}
\label{ch:ka-winning-people-over}
\hlchaptermodality{380}

\begin{chapteropening}
\textbf{By the end of this chapter:} \emph{I can say to persuade, to motivate, to drive home, an inducement and to attract.}

Complete the last lesson of chapter 380: I can say to persuade, to motivate, to drive home, an inducement and to attract.
\end{chapteropening}

\section[\texorpdfstring{\kn{ಮನವೊಲಿಸು} (manavolisu)}{ಮನವೊಲಿಸು (manavolisu)}]{\kn{ಮನವೊಲಿಸು} (manavolisu) — to persuade}

\label{lesson:KA-C380-manavolisu}



\begin{quote}
Before the new one: say the Kannada for to oppose, then the Kannada for mutual.
\end{quote}

\subsection*{You'll want to know: \kn{ಮನವೊಲಿಸು}}
\textbf{\kn{ಮನವೊಲಿಸು}} — \emph{manavolisu} — \textquotedblleft{}to persuade\textquotedblright{}.

\begin{quote}
\kn{ಅಮ್ಮನನ್ನು} \kn{ಮನವೊಲಿಸಲು} \kn{ನನಗೆ} \kn{ಒಂದು} \kn{ವಾರ} \kn{ಬೇಕಾಯಿತು}.
\end{quote}

\begin{quote}
\emph{ammanannu manavolisalu nanage ondu vāra bēkāyitu.}
\end{quote}

\begin{quote}
It took me a week to persuade Mother.
\end{quote}

Past (I …): \kn{ಮನವೊಲಿಸಿದೆ}, \emph{manavoliside}.

The person persuaded takes -\kn{ಅನ್ನು}.

\begin{cousinweb}
\textbf{\kn{ಮನವೊಲಿಸು}} (\emph{manavolisu}) is built from two parts: \emph{mana}, \textquotedblleft{}mind\textquotedblright{} (Sanskrit \emph{manas}, CDIAL 9822), and \emph{olisu}, \textquotedblleft{}to win over\textquotedblright{}, from native \emph{ol} / \emph{oli}, \textquotedblleft{}to be pleased, to favour\textquotedblright{} (DEDR 1006).
\end{cousinweb}

\begin{grammarlens}[title={What you've built}]
One more word for giving reasons and opinions.
\end{grammarlens}

\subsection*{Your turn}
\begin{itemize}
\raggedright
  \item \emph{Say it:} \emph{manavolisu}
  \item \emph{Say it:} \emph{manavolisu}, once more
  \item \emph{Say it:} \emph{ammanannu manavolisalu nanage ondu vāra bēkāyitu.} — It took me a week to persuade Mother.
  \item \emph{Recall:} say the Kannada for to oppose, then the Kannada for mutual, then say \emph{manavolisu} again
\end{itemize}

\begin{tcolorbox}[breakable,colback=teal!4,colframe=teal!35!black,title={Before you move on}]
What does \kn{ಮನವೊಲಿಸು} mean? (\textquotedblleft{}To persuade\textquotedblright{}.) Say it once more.
\end{tcolorbox}
\section[\texorpdfstring{\kn{ಪ್ರೇರೇಪಿಸು} (prērēpisu)}{ಪ್ರೇರೇಪಿಸು (prērēpisu)}]{\kn{ಪ್ರೇರೇಪಿಸು} (prērēpisu) — to motivate; to urge}

\label{lesson:KA-C380-prerepisu}



\begin{quote}
Before the new one: say the Kannada for mutual, then the Kannada for to persuade.
\end{quote}

\subsection*{You'll want to know: \kn{ಪ್ರೇರೇಪಿಸು}}
\textbf{\kn{ಪ್ರೇರೇಪಿಸು}} — \emph{prērēpisu} — \textquotedblleft{}to motivate; to urge\textquotedblright{}.

\begin{quote}
\kn{ಶಿಕ್ಷಕಿ} \kn{ನನ್ನನ್ನು} \kn{ಓದಲು} \kn{ಪ್ರೇರೇಪಿಸಿದರು}.
\end{quote}

\begin{quote}
\emph{śikṣaki nannannu ōdalu prērēpisidaru.}
\end{quote}

\begin{quote}
The teacher motivated me to study.
\end{quote}

Past (I …): \kn{ಪ್ರೇರೇಪಿಸಿದೆ}, \emph{prērēpiside}.

The person urged takes -\kn{ಅನ್ನು}, the aim the -\kn{ಅಲು} form: \kn{ಓದಲು}, to study.

\begin{cousinweb}
\textbf{\kn{ಪ್ರೇರೇಪಿಸು}} (\emph{prērēpisu}) is usually traced to Sanskrit \emph{preraṇa}, \textquotedblleft{}setting in motion, urging, inciting\textquotedblright{} (Monier-Williams), from \emph{pra-īr}, \textquotedblleft{}to urge on\textquotedblright{}, with the Kannada verb ending -\emph{isu}; how the shape -\emph{ēpisu} arose is not settled.
\end{cousinweb}

\begin{grammarlens}[title={What you've built}]
One more word for giving reasons and opinions.
\end{grammarlens}

\subsection*{Your turn}
\begin{itemize}
\raggedright
  \item \emph{Say it:} \emph{prērēpisu}
  \item \emph{Say it:} \emph{prērēpisu}, once more
  \item \emph{Say it:} \emph{śikṣaki nannannu ōdalu prērēpisidaru.} — The teacher motivated me to study.
  \item \emph{Recall:} say the Kannada for mutual, then the Kannada for to persuade, then say \emph{prērēpisu} again
\end{itemize}

\begin{tcolorbox}[breakable,colback=teal!4,colframe=teal!35!black,title={Before you move on}]
What does \kn{ಪ್ರೇರೇಪಿಸು} mean? (\textquotedblleft{}To motivate; to urge\textquotedblright{}.) Say it once more.
\end{tcolorbox}
\section[\texorpdfstring{\kn{ಮನದಟ್ಟು} \kn{ಮಾಡು} (manadaṭṭu māḍu)}{ಮನದಟ್ಟು ಮಾಡು (manadaṭṭu māḍu)}]{\kn{ಮನದಟ್ಟು} \kn{ಮಾಡು} (manadaṭṭu māḍu) — to drive home; to make clear}

\label{lesson:KA-C380-manadattu-madu}



\begin{quote}
Before the new one: say the Kannada for to persuade, then the Kannada for to motivate.
\end{quote}

\subsection*{You'll want to know: \kn{ಮನದಟ್ಟು} \kn{ಮಾಡು}}
\textbf{\kn{ಮನದಟ್ಟು} \kn{ಮಾಡು}} — \emph{manadaṭṭu māḍu} — \textquotedblleft{}to drive home; to make clear\textquotedblright{}.

\begin{quote}
\kn{ನಿಯಮದ} \kn{ಉದ್ದೇಶವನ್ನು} \kn{ಶಿಕ್ಷಕರು} \kn{ಮಕ್ಕಳಿಗೆ} \kn{ಮನದಟ್ಟು} \kn{ಮಾಡಿದರು}.
\end{quote}

\begin{quote}
\emph{niyamada uddēśavannu śikṣakaru makkaḷige manadaṭṭu māḍidaru.}
\end{quote}

\begin{quote}
The teacher made the purpose of the rule clear to the children.
\end{quote}

Past (I …): \kn{ಮನದಟ್ಟು} \kn{ಮಾಡಿದೆ}, \emph{manadaṭṭu māḍide}.

The listener takes -\kn{ಗೆ}. \kn{ಮನದಟ್ಟಾಗು}, to sink in.

\begin{cousinweb}
\textbf{\kn{ಮನದಟ್ಟು} \kn{ಮಾಡು}} (\emph{manadaṭṭu māḍu}) is built from \emph{mana}, \textquotedblleft{}mind\textquotedblright{} (Sanskrit \emph{manas}, CDIAL 9822), joined to \emph{daṭṭu}, whose history is not settled, with \emph{māḍu}, \textquotedblleft{}to do\textquotedblright{}.
\end{cousinweb}

\begin{grammarlens}[title={What you've built}]
One more word for giving reasons and opinions.
\end{grammarlens}

\subsection*{Your turn}
\begin{itemize}
\raggedright
  \item \emph{Say it:} \emph{manadaṭṭu māḍu}
  \item \emph{Say it:} \emph{manadaṭṭu māḍu}, once more
  \item \emph{Say it:} \emph{niyamada uddēśavannu śikṣakaru makkaḷige manadaṭṭu māḍidaru.} — The teacher made the purpose of the rule clear to the children.
  \item \emph{Recall:} say the Kannada for to persuade, then the Kannada for to motivate, then say \emph{manadaṭṭu māḍu} again
\end{itemize}

\begin{tcolorbox}[breakable,colback=teal!4,colframe=teal!35!black,title={Before you move on}]
What does \kn{ಮನದಟ್ಟು} \kn{ಮಾಡು} mean? (\textquotedblleft{}To drive home; to make clear\textquotedblright{}.) Say it once more.
\end{tcolorbox}
\section[\texorpdfstring{\kn{ಆಮಿಷ} (āmiṣa)}{ಆಮಿಷ (āmiṣa)}]{\kn{ಆಮಿಷ} (āmiṣa) — an inducement; a lure}

\label{lesson:KA-C380-amisa}



\begin{quote}
Before the new one: say the Kannada for to motivate, then the Kannada for to drive home.
\end{quote}

\subsection*{You'll want to know: \kn{ಆಮಿಷ}}
\textbf{\kn{ಆಮಿಷ}} — \emph{āmiṣa} — \textquotedblleft{}an inducement; a lure\textquotedblright{}.

\begin{quote}
\kn{ಹಣದ} \kn{ಆಮಿಷಕ್ಕೆ} \kn{ಅವನು} \kn{ಒಪ್ಪಲಿಲ್ಲ}.
\end{quote}

\begin{quote}
\emph{haṇada āmiṣakke avanu oppalilla.}
\end{quote}

\begin{quote}
He did not give in to the lure of money.
\end{quote}

\kn{ಆಮಿಷ} \kn{ಒಡ್ಡು}, to dangle an inducement.

\begin{cousinweb}
\textbf{\kn{ಆಮಿಷ}} (\emph{āmiṣa}) is borrowed from Sanskrit \emph{āmiṣa}, \textquotedblleft{}flesh, prey\textquotedblright{}, and so \textquotedblleft{}a bait\textquotedblright{} (CDIAL 1256).
\end{cousinweb}

\begin{grammarlens}[title={What you've built}]
One more word for giving reasons and opinions.
\end{grammarlens}

\subsection*{Your turn}
\begin{itemize}
\raggedright
  \item \emph{Say it:} \emph{āmiṣa}
  \item \emph{Say it:} \emph{āmiṣa}, once more
  \item \emph{Say it:} \emph{haṇada āmiṣakke avanu oppalilla.} — He did not give in to the lure of money.
  \item \emph{Recall:} say the Kannada for to motivate, then the Kannada for to drive home, then say \emph{āmiṣa} again
\end{itemize}

\begin{tcolorbox}[breakable,colback=teal!4,colframe=teal!35!black,title={Before you move on}]
What does \kn{ಆಮಿಷ} mean? (\textquotedblleft{}An inducement; a lure\textquotedblright{}.) Say it once more.
\end{tcolorbox}
\section[\texorpdfstring{\kn{ಆಕರ್ಷಿಸು} (ākarṣisu)}{ಆಕರ್ಷಿಸು (ākarṣisu)}]{\kn{ಆಕರ್ಷಿಸು} (ākarṣisu) — to attract}

\label{lesson:KA-C380-akarsisu}



\begin{quote}
Before the new one: say the Kannada for to drive home, then the Kannada for an inducement.
\end{quote}

\subsection*{You'll want to know: \kn{ಆಕರ್ಷಿಸು}}
\textbf{\kn{ಆಕರ್ಷಿಸು}} — \emph{ākarṣisu} — \textquotedblleft{}to attract\textquotedblright{}.

\begin{quote}
\kn{ಹೊಸ} \kn{ಉದ್ಯಾನ} \kn{ಅನೇಕ} \kn{ಜನರನ್ನು} \kn{ಆಕರ್ಷಿಸಿತು}.
\end{quote}

\begin{quote}
\emph{hosa udyāna anēka janarannu ākarṣisitu.}
\end{quote}

\begin{quote}
The new park attracted many people.
\end{quote}

Past (I …): \kn{ಆಕರ್ಷಿಸಿದೆ}, \emph{ākarṣiside}.

\kn{ಆಕರ್ಷಕ}, attractive; \kn{ಆಕರ್ಷಣೆ}, attraction.

\begin{cousinweb}
\textbf{\kn{ಆಕರ್ಷಿಸು}} (\emph{ākarṣisu}) is built on Sanskrit \emph{ākarṣaṇa}, \textquotedblleft{}pulling, drawing near, attracting\textquotedblright{} (Monier-Williams), with the Kannada verb ending -\emph{isu}.
\end{cousinweb}

\begin{grammarlens}[title={What you've built}]
One more word for giving reasons and opinions.
\end{grammarlens}

\subsection*{Your turn}
\begin{itemize}
\raggedright
  \item \emph{Say it:} \emph{ākarṣisu}
  \item \emph{Say it:} \emph{ākarṣisu}, once more
  \item \emph{Say it:} \emph{hosa udyāna anēka janarannu ākarṣisitu.} — The new park attracted many people.
  \item \emph{Recall:} say the Kannada for to drive home, then the Kannada for an inducement, then say \emph{ākarṣisu} again
\end{itemize}

\begin{tcolorbox}[breakable,colback=teal!4,colframe=teal!35!black,title={Before you move on}]
What does \kn{ಆಕರ್ಷಿಸು} mean? (\textquotedblleft{}To attract\textquotedblright{}.) Say it once more.
\end{tcolorbox}
